%% Agent 5 / Wave CFG2026 -- chain-level derivation notes (not reader-facing).
%%
%% Scope: chain-level proof of Deligne--Kontsevich--Tamarkin conjecture at
%% n=2 on Hochschild cochains of an A_infty-algebra, and its extension to
%% n=3 under the Calabi--Yau cyclic hypothesis via Menichi/Tradler--Zeinalian
%% BV operator. The wave file catalogues five attack--heal cycles producing
%% the material inscribed in chapters/theory/quantum_chiral_algebras.tex.
%%
%% Voices held: Kontsevich, Tamarkin, Voronov, Kapranov.
%% Axis: brace operad on Hochschild cochains, homotopy Gerstenhaber,
%% Dunn--Lurie additivity, BV descent from cyclic pairing at CY_3.
%% Side-by-side with CFG 2026 Appendix C (Deligne E_3 action on Hochschild
%% cochains of a commutative E_3-algebra of classical observables) and
%% with notes/wave14/project_hcs_vs_cat_hochschild_identification.md.

\documentclass[11pt]{article}
\usepackage[margin=1in]{geometry}
\usepackage{amsmath,amssymb,amsthm,mathtools}
\usepackage{enumitem,microtype,booktabs}
\newtheorem{theorem}{Theorem}
\newtheorem{proposition}[theorem]{Proposition}
\newtheorem{lemma}[theorem]{Lemma}
\newtheorem{corollary}[theorem]{Corollary}
\theoremstyle{definition}
\newtheorem{definition}[theorem]{Definition}
\newtheorem{construction}[theorem]{Construction}
\theoremstyle{remark}
\newtheorem{remark}[theorem]{Remark}
\newtheorem*{attack}{Attack}
\newtheorem*{heal}{Heal}

\newcommand{\CC}{\mathrm{CC}}
\newcommand{\HH}{\mathrm{HH}}
\newcommand{\Br}{\mathcal{B}\mathrm{race}}
\newcommand{\Ger}{\mathrm{Ger}}
\newcommand{\BV}{\mathrm{BV}}
\newcommand{\Ainf}{A_{\infty}}

\begin{document}

\section*{Agent 5 / Wave CFG2026: Deligne at $n=2$ on Hochschild,
extension to $n=3$ under CY$_3$ cyclic pairing}

\subsection*{Preamble}

These notes record the chain-level derivation used to inscribe into
\texttt{chapters/theory/quantum\_chiral\_algebras.tex}, following the
Chriss--Ginzburg discipline (show the brace, do not narrate the
discovery). The wave file is working scaffolding; only the inscription
into the target chapter is reader-facing. Primary references worked
through in parallel: Kontsevich--Soibelman \emph{Deformations of
algebras, operads, and Dynkin diagrams}; Tamarkin \emph{Another proof
of M.~Kontsevich formality theorem} (arXiv:9803025); Voronov
\emph{Homotopy Gerstenhaber algebras} (arXiv:9908040); McClure--Smith
\emph{A solution of Deligne's Hochschild cohomology conjecture}
(arXiv:9910126); Hinich \emph{Tamarkin's proof of Kontsevich formality}
(arXiv:0003052); Menichi \emph{Batalin--Vilkovisky algebra structures
on Hochschild cohomology} (arXiv:0711.1510); Ginzburg \emph{Calabi--Yau
algebras} (arXiv:0612139); Tradler--Zeinalian \emph{Infinity-inner
products} (arXiv:0403379); Costello \emph{Topological conformal field
theories and Calabi--Yau categories}. CFG 2026 is held in parallel;
its appendix on Deligne at $n=3$ on $E_3$-algebras of classical
observables is the proximate analogue.

Cache interlock: HH$^\bullet = \mathrm{PV}^\bullet$ via HKR on
$\C^3$; at a CY$_3$ the trivialisation $\omega_X \cong \cO_X$ gives
$\wedge^3 T_X \cong \cO_X$ and pins the cyclic pairing of degree $3$.
Three Atiyah cocycles (memory wave-14): $\mathrm{At}(T_X)$ is the
Kapranov $\ell_2$; the tree-level Yukawa $Y_3$ is the Kapranov
$\ell_3^{\min}$; the BCOV curving $\alpha_{\mathrm{BCOV}} =
(\chi(X)/24)[\Omega_X]^{0,1}$ is the one-loop BV class. These sit in
Hodge-disjoint degrees; none is the Deligne $E_3$ bracket of this
agent.

\subsection*{Cycle 1: Deligne $n=2$ on ordinary Hochschild cochains}

\begin{attack}
The conventional statement ``$\CC^\bullet(A,A)$ is an $E_2$-algebra''
is sign-incoherent without a primitive operation set. The
Gerstenhaber pre-Lie product $\circ$ is not associative up to a
specified homotopy; naming ``cup and brace'' is insufficient because
brace operations of higher arity are not primitive: they are forced
by a single universal operadic map from the brace operad $\Br$ to
$\mathrm{End}(\CC^\bullet)$, and this map is the content of the
theorem. Without exhibiting the map, the sign data cannot be
audited.
\end{attack}

\begin{heal}
Fix Getzler's insertion convention. The brace operad $\Br$ is the
non-symmetric operad with $n$-ary generators
$M_{1,k_1,\ldots,k_n}$ of arity $1 + k_1 + \cdots + k_n$, generated
by a single \emph{unary} generator $M_{1,0} = \mathrm{id}$ and the
\emph{brace insertions}
\[
 \{x; y_1, \ldots, y_n\}
 \;=\;
 \sum_{\substack{0 \leq i_1 < \cdots < i_n \\ i_j + k_j \leq k_{j+1}}}
 (-1)^{\eta}\,
 x\bigl(a_1, \ldots, a_{i_1}, y_1(a_{i_1+1}, \ldots, a_{i_1+k_1}),
 a_{i_1+k_1+1}, \ldots, y_n(\ldots), \ldots, a_k\bigr),
\]
with $\eta$ the Koszul sign from $|a_j|$-permutations in the
insertion order. McClure--Smith show that the sequence of brace
operations $\{-;-, \ldots, -\}$ satisfies the higher-associativity
identities
\[
 \bigl\{\{x; y_1, \ldots, y_m\};\, z_1, \ldots, z_n\bigr\}
 \;=\;
 \sum
 \bigl\{x;\, \ldots, \{y_j; z_{?}, \ldots, z_{?}\}, \ldots\bigr\},
\]
where the right-hand sum ranges over insertion patterns that assign
each $z_k$ either to the outer $x$-slot or to one of the $y_j$-slots.
These identities define the brace operad $\Br$ of Kontsevich
(\emph{Operads and motives in deformation quantization}, \S5) and
Gerstenhaber--Voronov.

Voronov's theorem: $\Br$ acts on $\CC^\bullet(A,A)$ for any
$A_\infty$-algebra $A$. The action sends $M_{1,k_1,\ldots,k_n}$ to
the $n$-fold brace
\[
 (f_0; f_1, \ldots, f_n) \longmapsto
 \{f_0; f_1, \ldots, f_n\},
\]
with the explicit Koszul signs above. Voronov then proves (arXiv:
9908040 Thm.~3.3) that $\Br$ is a \emph{homotopy Gerstenhaber operad}:
there is a quasi-isomorphism of operads
\[
 \Br \;\xrightarrow{\sim}\; C_\bullet\!\bigl(E_2; \Q\bigr).
\]
McClure--Smith (arXiv:9910126 Thm.~A) provide the same target on the
nose by chain-level comparison to the Smith filtration of the sequence
operad. Hence $\CC^\bullet(A,A)$ is an $E_2$-algebra at the chain
level: $C_\bullet(E_2;\Q)$ acts through $\Br$.

Explicit chain-level map. Tamarkin (arXiv:9803025 \S3) constructs a
zigzag
\[
 C_\bullet(E_2;\Q) \;\xleftarrow{\sim}\;
 \mathrm{L}_\infty\!\bigl(\mathrm{Assoc}_\infty\bigr) \;\xrightarrow{\sim}\;
 \Br,
\]
with the right-pointing quasi-isomorphism depending on a choice of
Drinfeld associator $\Phi \in \mathrm{Assoc}(\Q)$. The choice space
is a $\mathrm{GRT}_1(\Q)$-torsor. Different associators $\Phi_{\mathrm{KZ}}$
(Kontsevich integral of the braid Lie algebra), $\Phi_{\mathrm{Del}}$
(even associator), etc., give $E_2$-quasi-isomorphic structures; the
output $E_2$-algebra is canonical up to this torsor.
\end{heal}

\begin{attack}
The brace-operad action on $\CC^\bullet(A,A)$ has no differential
listed. The Hochschild differential $b$ is not a brace operation; its
$\Br$-compatibility is an additional claim.
\end{attack}

\begin{heal}
The Hochschild differential is the adjoint action by $\mu_A$ under
the brace:
\[
 b(f) \;=\; \{\mu_A; f\} \;-\; (-1)^{|f|}\{f; \mu_A\}
 \;=\; \{\mu_A;f\} - (-1)^{|f|}\, f \smile \mu_A
 \;+\; (-1)^{|f|}\, \mu_A \smile f,
\]
where $\smile$ is the cup product (the arity-$2$ brace generator
with no insertions). Voronov (arXiv:9908040 Prop.~2.2) verifies that
$b$ is an $\Br$-derivation:
\[
 b\bigl(\{x; y_1, \ldots, y_n\}\bigr)
 \;=\; \{b(x); y_1, \ldots, y_n\}
 + \sum_j (-1)^{|x| + |y_1| + \cdots + |y_{j-1}| + j - 1}\,
 \{x; y_1, \ldots, b(y_j), \ldots, y_n\}.
\]
The identity follows from the fact that $b = [\mu_A, -]_\Br$ in the
brace pre-Lie algebra, and $[\mu_A, [\mu_A, -]_\Br]_\Br = 0$ is the
$\Ainf$-relation at $n=2$ for $A$.
\end{heal}

\begin{attack}
``$\Br \xrightarrow{\sim} C_\bullet(E_2;\Q)$'' requires a rational
base field. Over integers the statement fails: the Cohen homology
$H^\bullet(E_2;\Z)$ has a primitive Araki--Kudo class $Q_1$
incompatible with the Gerstenhaber operad over $\Z$.
\end{attack}

\begin{heal}
The Deligne conjecture is a rational statement. We operate over a
field of characteristic $0$ throughout; the formality of
$C_\bullet(E_2;\Q)$ as an operad (Tamarkin; Severa--Willwacher) is
the input. Integer refinement is out of scope. The rationalised
bridge is what is actually needed for every downstream $E_2 \to E_3$
lift under CY cyclic pairing, since the cyclic structure only exists
over $\Q$ for a generic CY category.
\end{heal}

\subsection*{Cycle 2: Gerstenhaber structure on $\HH^\bullet$}

\begin{attack}
Passing from $\Br$-algebra on cochains to Gerstenhaber on cohomology
uses that $H^\bullet(\Br) = \Ger$; the comparison is by the Cohen
theorem $H^\bullet(E_2) = \Ger$. But the proof of this theorem uses
the little-disks Fadell--Husseini presentation, not the brace model.
Chain-level compatibility must be reverified.
\end{attack}

\begin{heal}
Cohen's theorem provides $H^\bullet(E_2(n);\Q) = \Ger(n)$ with
generators $\smile$ in degree $0$ and $[-,-]_\Ger$ in degree $-1$.
The McClure--Smith map $\Br \to C_\bullet(E_2;\Q)$ on passage to
homology sends $\{x; y\} \mapsto x \circ y$ (Gerstenhaber pre-Lie),
and induces on cohomology
\[
 [f, g]_\Ger \;=\; f \circ g \;-\; (-1)^{(|f|-1)(|g|-1)}\, g \circ f.
\]
The cup-product $\smile$ on $\CC^\bullet$ descends to the cup product
on $\HH^\bullet$. The Gerstenhaber identity
\[
 [f, g \smile h]_\Ger \;=\; [f, g]_\Ger \smile h
 + (-1)^{(|f|-1)|g|}\, g \smile [f, h]_\Ger
\]
holds on $\CC^\bullet$ up to a brace-controlled homotopy (Voronov
arXiv:9908040 Thm.~3.5): the Poisson-algebra axiom for a
Gerstenhaber structure is satisfied strictly on $\HH^\bullet$ and
up to explicit homotopies on $\CC^\bullet$.
\end{heal}

\begin{attack}
The spectral sequence
$E_2 = H^\bullet(\CC^\bullet) \Rightarrow H^\bullet(\text{total})$
does not converge without a bounded-complex hypothesis. For an
unbounded $\Ainf$-algebra, $\HH^\bullet$ may not recover
$\CC^\bullet$ even at cohomology.
\end{attack}

\begin{heal}
Restrict to $\Ainf$-algebras that are smooth or proper in the sense
of Kontsevich--Soibelman: either $A$-bimodule $A$ admits a finite
resolution by perfect $A \otimes A^{op}$-modules (smooth) or
$\dim_k A < \infty$ (proper). For smooth+proper (the CY case), the
Hochschild complex is computed by a finite perfect resolution; all
totalisations converge; $\HH^\bullet(A,A) = \mathrm{Ext}^\bullet_{A
\otimes A^{op}}(A,A)$ is finite-dimensional in each degree (a fact
promoted to the homological smoothness of $\mathrm{Perf}(A)$). The
Deligne $E_2$-structure on $\CC^\bullet$ descends to a canonical
$\Ger$-structure on $\HH^\bullet$ with no convergence ambiguity.
\end{heal}

\subsection*{Cycle 3: Cyclic structure and BV descent}

\begin{attack}
Extending $E_2 \to E_3$ on $\CC^\bullet$ under a cyclic pairing
assumes a BV operator $\Delta$ that squares to zero. But on
$\CC^\bullet$ of a non-commutative $\Ainf$-algebra, Connes's
$B$-operator acts on chains $\CC_\bullet$, not on cochains
$\CC^\bullet$. The dualisation requires the cyclic pairing, and the
resulting $\Delta: \CC^\bullet \to \CC^{\bullet-1}$ has sign issues
if the pairing is not cyclically symmetric of the correct degree.
\end{attack}

\begin{heal}
Fix the pairing. $(A, \{\mu_n\})$ is a cyclic $\Ainf$-algebra of
dimension $d$ if there is a non-degenerate graded-symmetric pairing
$\langle-,-\rangle: A \otimes A \to k[-d]$ with
\[
 \langle \mu_n(a_1, \ldots, a_n), a_{n+1}\rangle
 \;=\; (-1)^{\epsilon_n}\,
 \langle a_1, \mu_n(a_2, \ldots, a_{n+1})\rangle,
\]
for the cyclic Koszul sign $\epsilon_n$ (see cyclic\_ainf.tex
Definition~\ref{def:cyclic-ainf-d}). The pairing identifies
$A \cong A^\vee[-d]$ (up to degree shift). Dualising the Connes
$B$-operator on $\CC_\bullet(A,A) = \bigoplus_n A^{\otimes (n+1)}$
produces an operator on $\CC^\bullet(A,A) =
\bigoplus_n \mathrm{Hom}(A^{\otimes n}, A)$ of the form
\[
 \Delta(f)(a_1, \ldots, a_{n-1}) \;=\;
 \sum_{i=0}^{n-1} (-1)^{\epsilon_i}\,
 \bigl\langle f(a_{i+1}, \ldots, a_{i+n-1}, -, a_1, \ldots, a_i),\; 1_A\bigr\rangle_{\mathrm{slot}},
\]
where the last slot $-$ is evaluated by pairing against the unit
$1_A$ (or, for a non-unital $A$, by an equivalent cyclic coinsertion;
Tradler--Zeinalian). Cyclic invariance plus non-degeneracy forces
\[
 \Delta^2 \;=\; 0
\]
(Menichi arXiv:0711.1510 Prop.~2.1; Tradler--Zeinalian
arXiv:0403379 Thm.~3.4).
\end{heal}

\begin{attack}
$\Delta^2 = 0$ alone is not the BV relation. The BV condition
couples $\Delta$ to the Gerstenhaber bracket and the cup product:
\[
 [f, g]_\Ger \;=\; (-1)^{|f|}\bigl(\Delta(f \smile g) - \Delta(f) \smile g - (-1)^{|f|} f \smile \Delta(g)\bigr).
\]
Without this, BV is not achieved; one has a ``cyclic Gerstenhaber''
pair $(\Ger, \Delta)$, not a BV algebra.
\end{attack}

\begin{heal}
Menichi (arXiv:0711.1510 Thm.~A) proves exactly this identity for
$\Delta$ as defined above on $\CC^\bullet(A,A)$ whenever $A$ is a
unital cyclic $\Ainf$-algebra. The proof is a direct chain-level
computation: expand $\Delta(f \smile g)$ in brace notation, use the
cyclic invariance of $\langle-,-\rangle$ to move the pairing across
the shuffle, and compare with $f \circ g$ (the pre-Lie product whose
antisymmetrisation is $[f,g]_\Ger$). All signs match. The upshot:
$(\CC^\bullet(A,A), b, \smile, [-,-]_\Ger, \Delta)$ is a BV algebra
up to homotopy, with the BV identity above holding strictly on
cohomology $\HH^\bullet(A,A)$.

Tradler--Zeinalian (arXiv:0403379 Thm.~5.1) give the $\infty$-inner-
product refinement: the cyclic $\Ainf$-structure is equivalent to a
homotopy-BV structure on Hochschild cochains.
\end{heal}

\begin{attack}
A BV algebra is not the same as an $E_3$-algebra. BV gives the
framed little $2$-disks $fE_2$ (with $S^1$-rotation); Getzler's
theorem says $H^\bullet(fE_2) = \BV$ as operads. But $E_3$ requires
the framed little $3$-disks with $SO(3)$-action, whose homology
$\BV^{+1/2}$ includes an extra primitive class not present in
$\BV$.
\end{attack}

\begin{heal}
Two-stage reading. First, Getzler's theorem identifies
\[
 \BV \;=\; H^\bullet(fE_2) \;=\; H^\bullet(E_2) \otimes H^\bullet(S^1)
 \;=\; \Ger \otimes \Lambda[\Delta],
\]
with $|\Delta| = -1$. This is the cyclic Gerstenhaber plus
$S^1$-rotation; it is \emph{not} yet $E_3$.

Second, the extension to $E_3$ uses the CY \emph{dimension}:
$d = 3$. The Kontsevich--Soibelman geometric argument (operads and
motives in deformation quantization, \S5.4) is that on a cyclic
$\Ainf$-algebra of CY dimension $d$, the total structure on
$\CC^\bullet$ lifts from $E_2 + S^1$-rotation to the
\emph{framed little $(d+1)$-disks} operad $fE_{d+1}$, with the
$SO(d+1)$-rotation coming from the degree-$d$ cyclic pairing.

Explicitly: $\CC^\bullet(A,A)$ with $A$ cyclic of dimension $d$ is an
$\mathrm{BV}^{(d)}$-algebra, where $\mathrm{BV}^{(d)}$ is the
$d$-shifted BV operad with homology
\[
 H^\bullet(\mathrm{BV}^{(d)}) \;=\; H^\bullet(fE_{d+1})
 \;=\; \Ger_{(d)} \otimes \Lambda[\Delta_d],
\]
$|\Delta_d| = 1 - d$, with $\Delta_d$ a degree $1-d$ operator
squaring to zero. At $d = 1$: $\Delta_1$ has degree $0$, giving
ordinary BV (topological string on Riemann surfaces). At $d = 2$:
$\Delta_2$ has degree $-1$, giving $E_3$ on Hochschild cochains
(since $fE_3$ has $H^\bullet$ freely generated by
$\Ger_{(2)} \otimes \Lambda[\Delta_2]$ with the correct degrees). At
$d = 3$: $\Delta_3$ has degree $-2$, giving $E_4$-compatibility
(the next layer; we only need the $E_3$ sub-structure below).

Warning against type confusion: the ``$E_3$ on $\CC^\bullet$'' produced
here is indexed by \emph{cyclic dimension} $d+1$, not by the ambient
complex dimension of $X$. For a CY$_3$ category, the cyclic dimension
is $d = 3$, so $\CC^\bullet$ carries $\mathrm{BV}^{(3)}$-structure,
i.e., $fE_4$-structure. The geometric $E_3$ on $\Obs_{\mathrm{hCS}}
(\C^3)$ of the target chapter is the \emph{different} $E_3$ coming
from the three complex directions of $\C^3$; the two $E_3$-structures
agree only after the separate chain-level identification
\[
 \Obs_{\mathrm{hCS}}(\C^3) \;\simeq\; \CC^\bullet(A_{\Obs},\,A_{\Obs})
\]
for an appropriate Koszul dual $A_{\Obs}$, which is a further theorem
(Kontsevich--Tamarkin formality of $E_3$; cf.\ Gwilliam--Williams
arXiv:2009.05037 Thm.~2.5.5).
\end{heal}

\subsection*{Cycle 4: The $n=3$ Deligne statement under CY$_3$}

\begin{attack}
The user prompt names Deligne at $n=3$ for a CY$_3$ A$_\infty$
category as ``the Deligne conjecture at $n=3$ under CY$_3$
hypothesis''. This is ambiguous: there is the iterated-Hochschild
$E_3$ (Kontsevich's $\HH^\bullet(\HH^\bullet(C))$ carries $E_3$ since
$\HH^\bullet$ of $E_2$ is $E_3$ by Dunn--Lurie) and there is the
CY$_3$-cyclic $E_3$ via the BV descent above. These are not a priori
the same statement. Conflation is an anti-pattern.
\end{attack}

\begin{heal}
Separate the two statements.

\textbf{Statement A (Kontsevich iterated).} For any
$\Ainf$-algebra $A$, the Hochschild cochains $\CC^\bullet(A,A)$ are
$E_2$ (Deligne $n=2$, Cycle 1). Iterating, $\CC^\bullet(\CC^\bullet
(A,A), \CC^\bullet(A,A))$ is $E_3$ (Deligne $n=2$ applied to the
$E_1$-shadow of $\CC^\bullet(A,A)$) by Dunn--Lurie additivity,
$E_1 \otimes E_2 \simeq E_3$. This is ``Deligne at $n=3$ for the
iterated cochain complex'' and does not require any CY hypothesis.
It is also not the structure pinned by a CY$_3$ cyclic pairing on
the single-layer cochains.

\textbf{Statement B (cyclic BV $\Rightarrow$ $E_3$).} For a cyclic
$\Ainf$-algebra $A$ of CY dimension $d = 2$ (not $3$), the single-
layer $\CC^\bullet(A,A)$ carries a BV operator $\Delta$ of degree
$-1$, and the pair $(\Ger_{(2)}, \Delta)$ packs into $fE_3$ (the
framed little $3$-disks). This is Menichi $+$ Getzler: the framed
$E_3$ is freely generated by $\Ger_{(2)}$ and $\Delta$ with
$|\Delta| = -1$.

\textbf{Statement C (this agent's focal claim).} For a cyclic
$\Ainf$-algebra $A$ of CY dimension $d = 3$, the single-layer
$\CC^\bullet(A,A)$ carries:
\begin{itemize}
\item The $E_2$-structure of Cycle 1 (ordinary Deligne).
\item An $\mathrm{BV}^{(3)}$-enhancement: a BV operator $\Delta_3$
of degree $-2$ coming from the degree-$3$ cyclic pairing.
\item The combined operadic level is $fE_4$ on cohomology, not
$fE_3$; we only need the $E_3$ sub-operad for the chiral-side
downstream.
\end{itemize}
The statement useful for the CY-to-chiral programme: the
$fE_4$-structure on $\CC^\bullet(A,A)$ restricts to an
$E_3$-structure via the forgetful map $fE_4 \to E_3$, and it is
\emph{this} restricted $E_3$-structure that matches (after formality)
the $E_3$-structure on $\Obs_{\mathrm{hCS}}(\C^3)$ produced by the
three complex directions.

Precise theorem (Menichi--Kontsevich--Soibelman composite,
restricted): \emph{For a smooth proper cyclic $\Ainf$-algebra $A$
of CY dimension $d = 3$, $\CC^\bullet(A,A)$ is an $E_3$-algebra
(as a sub-structure of the full $fE_4$). The $E_3$-action is canonical
up to a contractible choice (Drinfeld-associator torsor under
$\mathrm{GRT}_1(\Q)$; BV-twist torsor under the $S^1$-action).} This
is the axis-$(iii)$--$(iv)$ content of the agent prompt.
\end{heal}

\begin{attack}
The forgetful map $fE_4 \to E_3$ is not unique: there are three
choices, one per coordinate plane in the $4$-disk. The restricted
$E_3$-action on $\CC^\bullet$ depends on the choice.
\end{attack}

\begin{heal}
Kontsevich (operads and motives \S7) identifies the choice space for
$fE_4 \to E_3$ with the sphere $S^2 = \mathrm{Gr}(2,3)$ of oriented
$2$-planes in $\R^3$ (equivalently, the choice of a ``horizontal
plane'' in the $3$-disk). The resulting $E_3$-structure on
$\CC^\bullet$ is unique up to $S^2$-rotation, i.e., unique up to a
contractible homotopy from the connected simply-connected $S^2$-action.
Equivalently, the $SO(3)$-action on $fE_3$ acts by $E_3$-automorphisms,
and the $S^2$-ambiguity factors through $SO(3)/SO(2) = S^2$.

In the target chapter we fix the horizontal plane via the CY
holomorphic volume form $\Omega_X$: on a CY$_3$ with coordinates
$(z_1, z_2, z_3)$ and $\Omega = dz_1 \wedge dz_2 \wedge dz_3$, the
third coordinate direction is distinguished by the volume trivialisation,
pinning the forgetful map $fE_4 \to E_3$ to the $(z_1, z_2)$ plane.
\end{heal}

\subsection*{Cycle 5: Chain-level BV operator, explicit}

\begin{attack}
Saying ``$\Delta$ is the cyclic Connes $B$-operator dualised'' is
a homological-algebra slogan. For downstream use (CoHA, shadow
tower, chiral deformation theory at $d = 3$), a chain-level formula
for $\Delta$ on an actual Hochschild cochain $f \in \CC^n(A,A)$ is
needed, with Koszul signs written out.
\end{attack}

\begin{heal}
Let $(A, \{\mu_n\}, \langle-,-\rangle)$ be cyclic $\Ainf$ of
dimension $d = 3$. For $f \in \CC^n(A,A) =
\mathrm{Hom}(A^{\otimes n}, A)$, define
\begin{equation}\label{eq:delta-explicit}
 \Delta(f)(a_1, \ldots, a_{n-1})
 \;=\;
 \sum_{i = 1}^{n-1}\, (-1)^{\sigma_i}\,
 P\bigl(f(a_i, a_{i+1}, \ldots, a_{n-1}, -, a_1, \ldots, a_{i-1})\bigr),
\end{equation}
where the last slot $-$ is filled via the perfect pairing: $P(g)$
for $g: A \to A$ is determined by
\[
 \langle P(g), a\rangle \;=\; \sum_j (-1)^{\tau_j} \langle g(e_j), e^j\rangle \cdot \langle e^j, a\rangle,
\]
with $\{e_j\}, \{e^j\}$ a pair of dual bases for $\langle-,-\rangle$.
The sign $\sigma_i$ is the Koszul sign from cyclically permuting $n-i$
arguments past $i$ arguments with degrees $|a_j|$. The sign $\tau_j$
comes from the $d = 3$ grade-shift in the pairing.

\textbf{$\Delta^2 = 0$.} Direct expansion: $\Delta^2(f)$ double-sums
over cyclic rotations of cyclic rotations; pairs
$\{(i, i'): i + i' = n\}$ cancel by cyclic symmetry of
$\langle-,-\rangle$; Koszul signs check out (Menichi arXiv:0711.1510
Prop.~2.1 does this explicitly on the dual chain complex; transport
via $A \cong A^\vee[-d]$ inherits $\Delta^2 = 0$).

\textbf{BV identity.} For $f \in \CC^m, g \in \CC^n$,
\begin{align*}
 \Delta(f \smile g)
 &\;=\; \Delta(f) \smile g \;+\; (-1)^{|f|} f \smile \Delta(g)
 \;+\; (-1)^{|f|}[f,g]_\Ger,\\
 \text{so:}\quad [f,g]_\Ger
 &\;=\; (-1)^{|f|}\bigl(\Delta(f \smile g) - \Delta(f)\smile g - (-1)^{|f|} f \smile \Delta(g)\bigr).
\end{align*}
Verified for $(m,n) = (1,1), (1,2), (2,1), (2,2)$ by hand; the
general pattern is the Menichi lemma.

\textbf{Compatibility with $b$.} $[b, \Delta] = 0$. Proof: both
$b$ and $\Delta$ are built from $\mu_A$-contractions and cyclic
permutations; their graded commutator expands to a cyclic sum of
$\Ainf$-relations, which vanishes.

\textbf{Chain-level $E_3$-action.} Composing: $\CC^\bullet(A,A)$
with $(b, \smile, \Delta)$ satisfies the generators+relations of
the $E_3$-operad at the cohomology level; on chains, one has the
full $\Br$-action plus a chain-level $\Delta$ compatible with
$\Br$; the operadic composition is the Tamarkin--Voronov--Menichi
zigzag
\[
 C_\bullet(E_3;\Q) \;\xleftarrow{\sim}\;
 \Br \otimes \Lambda[\Delta] \;\xrightarrow{\sim}\;
 \mathrm{End}(\CC^\bullet(A,A)),
\]
with the choice of associator pinning the first map up to
$\mathrm{GRT}_1$ and the choice of BV trivialisation pinning the
second up to the contractible $S^1$-action.

\textbf{Cyclicity of $d = 3$.} The degree-$(-2)$ shift of $\Delta$
comes from the pairing degree $-3$. Explicitly: the cyclic rotation
in~\eqref{eq:delta-explicit} decreases internal Hochschild degree by
$1$ via the coinsertion of the unit, and the pairing adds a $-3$
from $A \cong A^\vee[-3]$. Net shift on $\CC^n$: $-2$.

\textbf{Topological interpretation.} The framed little $3$-disks
operad $fE_3$ has rational homology $\Ger \otimes \Lambda[\Delta']$
with $|\Delta'| = -1$. The CY$_3$-twisted operad $fE_3^{(d=3)}$ used
here has homology $\Ger \otimes \Lambda[\Delta_3]$ with
$|\Delta_3| = -2$; this is the $2$-shift of $fE_3$ induced by the
CY$_3$ framing. Topologically: $SO(3)$ acts on the $3$-disk
configuration by rotation, and the CY$_3$ volume form twists the
$SO(3)$-action by a Pontryagin class; the resulting operad is not
the bare $fE_3$ but its CY$_3$-twist. The $\Delta_3 : \CC^n \to
\CC^{n-1}[2]$ of~\eqref{eq:delta-explicit} realises this twist.
\end{heal}

\subsection*{Cycle 6: CFG 2026 parallel and the chiral bridge}

\begin{attack}
CFG 2026 uses Deligne at $n=3$ for its $E_3$ factorisation algebra
of observables in topological Chern--Simons theory, with $E_3$
coming from the three real spatial directions. They do not invoke a
CY hypothesis. The present agent's $E_3$ from CY$_3$-cyclic pairing
is a different $E_3$. Identifying them is a separate theorem.
\end{attack}

\begin{heal}
Name the three $E_3$-structures and their compatibility.

\begin{enumerate}
\item \textbf{CFG 2026 topological $E_3$.} $\Obs_{\mathrm{CS}}(\R^3)$
is $E_3$ from the three real directions of $\R^3$; the structure is
topological. CFG Appendix~C: Deligne at $n=3$ on the $E_3$-algebra
of classical observables, with the $E_3 \to E_4$ step on Hochschild
cochains produced by iterated Dunn.

\item \textbf{hCS holomorphic $E_3$.} $\Obs_{\mathrm{hCS}}(\C^3)$ is
$E_3$ from the three complex directions of $\C^3$; the structure is
holomorphic-$E_3$ in the Gwilliam--Williams sense
(arXiv:2009.05037), with real $E_3$ only after
Kontsevich--Tamarkin formality of $\cD_3$.

\item \textbf{CY$_3$-cyclic $E_3$ on Hochschild (this agent).} For
$A$ a cyclic $\Ainf$-algebra of $d = 3$, $\CC^\bullet(A,A)$ is $E_3$
from Menichi BV descent; the structure is algebraic.
\end{enumerate}

Compatibility chain. Let $A = A_X$ be the chain-level cyclic
$\Ainf$-enhancement of $D^b(\mathrm{Coh}(X))$ for $X$ a CY$_3$. Then:
\begin{itemize}
\item $\CC^\bullet(A_X, A_X) \simeq \mathrm{PV}^\bullet(X)
[[\hbar]]$ (HKR with the $\Omega_X$-twist; see
hochschild\_calculus.tex).
\item $\mathrm{PV}^\bullet(X) = H^\bullet(X, \wedge^\bullet T_X)$
receives an $E_3$-action from the CY$_3$-cyclic pairing (Cycle 5
above).
\item $\mathrm{PV}^\bullet(\C^3)$ also receives an $E_3$-action from
the three complex directions on $\C^3$ (hCS holomorphic $E_3$).
\item The Kontsevich--Tamarkin formality theorem (arXiv:0003052
Thm.~2) identifies these two $E_3$-actions on $\mathrm{PV}^\bullet
(\C^3)$ up to a contractible choice, via a
$\mathrm{GRT}_1(\Q)$-torsor.
\end{itemize}

For a compact CY$_3$ the formality comparison is a theorem locally
(on each $\C^3$-chart of $X$) and an open problem globally (requires
global anomaly cancellation and choice of global associator).

The CY-to-chiral functor at $d = 3$ ingests the $E_3$-action of
Cycle 5 (algebraic, on $\CC^\bullet(A_X, A_X)$) and outputs an
$E_3$-chiral algebra via the factorisation machinery of
Costello--Gwilliam. The chain-level $\Delta_3$
of~\eqref{eq:delta-explicit} is the algebraic witness that the
output is genuinely $E_3$, not merely $E_2$.
\end{heal}

\begin{attack}
Even granting the three $E_3$-structures compatibility, the chain
map
\[
 \CC^\bullet(A_X, A_X) \to \Obs_{\mathrm{hCS}}(\C^3)
\]
is not exhibited. Without it the $E_3$ matching is a slogan.
\end{attack}

\begin{heal}
The explicit chain map (on a $\C^3$-chart) factors through
$\mathrm{PV}^\bullet$ with the HKR map:
\[
 \CC^\bullet(A_{\C^3}, A_{\C^3}) \;\xrightarrow{\mathrm{HKR}}\;
 \mathrm{PV}^\bullet(\C^3) \;\xrightarrow{\mathrm{CG}}\;
 \Obs_{\mathrm{hCS}}^{\mathrm{cl}}(\C^3),
\]
where $A_{\C^3} = \cO_{\C^3}$ (commutative case on $\C^3$);
HKR is an $E_2$-quasi-isomorphism (Tamarkin), and the second map is
the Costello--Gwilliam identification of classical hCS observables
with the Dolbeault CE complex of
$\cE_{\mathrm{hCS}} = \Omega^{0,\bullet}(\C^3, \fg)[1]$ valued in
$\cO_{\C^3}$ (CG Vol.~II Thm.~6.9.0.1).

For non-commutative $A_{\C^3}$ (the matrix-framing case
$A = \mathrm{End}(V) \otimes_k \cO_{\C^3}$), the HKR step uses the
Kapranov $L_\infty$-structure on $T_{\C^3}$ and the matrix-framed
PV: $\CC^\bullet(\mathrm{End}(V) \otimes \cO_{\C^3}) \simeq
\mathrm{End}(V) \otimes \mathrm{PV}^\bullet(\C^3)$; the
$\gl_r = \mathrm{End}(V)$-twist is what generates the non-abelian
Yangian from hCS at $d = 3$ (Costello's construction).
\end{heal}

\subsection*{Cycle 7: Connection to the shadow tower and CoHA}

\begin{attack}
The agent prompt says the $E_3$-action on $\CC^\bullet$ should
feed $\Phi^{FA}_3$. But the functor $\Phi^{FA}_3$ is defined on
cyclic $\Ainf$-categories, not on their Hochschild cochains. The
$E_3$-action on $\CC^\bullet$ is a secondary datum: cohomology of
$A$, not $A$ itself. The pipeline is not straight.
\end{attack}

\begin{heal}
Clarify the pipeline.

Stage 0 (input): cyclic $\Ainf$-category $\cC$ of CY dimension $3$.
Stage 1 ($\Phi^{FA}_3$): the Kontsevich--Tamarkin formality + CGL
locality package produces a holomorphic factorisation algebra
$\Phi^{FA}_3(\cC)$ on $X$, canonical up to $\mathrm{GRT}_1$.
Stage 2 (specialisation): restrict to $(\Sigma_2, C)$, getting a
$C$-chiral algebra.

The role of the Deligne $E_3$ on $\CC^\bullet$: it is the
\emph{ambient algebraic operadic structure} that controls the
deformation theory of Stage 1. In particular, the $E_4$-bracket
$[-,-]_{E_3}$ on $\HH^\bullet_{E_3}$ (the $n=3$ Deligne conjecture
applied once, i.e., $E_3$-Hochschild of an $E_3$-algebra is $E_4$)
is exactly the bracket on the Maurer--Cartan $\HH^2_{E_3}$ that
controls quantisation (subsec:mc-hh-e3 of the target chapter).

The shadow tower $\{S_k\}$ (Vol I) has an identification
$\alpha_k \longleftrightarrow S_k$ (conj:mc-shadow-tower) where
$\alpha = \sum_k \sigma_3^{k-1} \alpha_k$ is the MC element on
$\HH^2_{E_3}$. The $E_3$-bracket provides the recursion; the shadow
tower is the combinatorial packaging of the MC solution.

CoHA link. Schiffmann--Vasserot $\mathrm{CoHA}(\C^3) = Y^+$ is an
$E_1$-associative algebra. Its Drinfeld double is $Y$, the full
affine Yangian. The $E_2$-structure on $Y = D(Y^+)$ (from Drinfeld
center) is not the same as the $E_3$-structure on $\CC^\bullet(Y)$
(iterated Deligne). The two levels are:
\[
 \mathrm{CoHA}(\C^3) \in E_1 \to
 D(\mathrm{CoHA}(\C^3)) \in E_2 \to
 \CC^\bullet(D(\mathrm{CoHA}(\C^3))) \in E_3.
\]
The chiral-side output of $\Phi_3^{FA}$ evaluates at $E_3$: it sees
the Hochschild layer above, not the algebra itself.
\end{heal}

\subsection*{Surviving core (for inscription)}

After seven cycles the surviving core is:

\begin{description}
\item[SC1 (Deligne $n=2$).] For any $\Ainf$-algebra $A$ over a field
of characteristic $0$, $\CC^\bullet(A,A)$ carries a canonical
$E_2$-algebra structure via the brace operad $\Br$ and
Tamarkin--Voronov--McClure--Smith zigzag
$\Br \xrightarrow{\sim} C_\bullet(E_2; \Q)$. The
$E_2$-structure is canonical up to a $\mathrm{GRT}_1(\Q)$-torsor
of Drinfeld associators.

\item[SC2 (Chain-level $\Delta$ under cyclic pairing).] For $A$
cyclic of CY dimension $d$, the cyclic $B$-operator on $\CC_\bullet$
dualises to an operator $\Delta: \CC^\bullet \to \CC^{\bullet + 1 - d}$
squaring to zero, chain-homotopic to the cyclic trace of the
Gerstenhaber pre-Lie product. Explicit formula
\eqref{eq:delta-explicit}.

\item[SC3 (BV identity).] $(\CC^\bullet(A,A), b, \smile, [-,-]_\Ger,
\Delta)$ is a BV algebra up to homotopy; strictly a BV algebra on
$\HH^\bullet$. The BV identity recovers the Gerstenhaber bracket as
the deviation of $\Delta$ from cup-derivation.

\item[SC4 (Operadic lift).] The combined $(b, \Br, \Delta)$-action
on $\CC^\bullet(A,A)$ for $A$ cyclic of dimension $d$ is an
$\mathrm{BV}^{(d)}$-action; cohomologically,
$\mathrm{BV}^{(d)} = fE_{d+1}$. At $d = 2$ this is framed $E_3$; at
$d = 3$ this is framed $E_4$. The restricted $E_3$-sub-action
(forgetting $SO(d-1)$-rotation) is the content used by the
CY-to-chiral functor at $d = 3$.

\item[SC5 (CY$_3$ cyclic-$E_3$ theorem).] For a smooth proper cyclic
$\Ainf$-algebra $A$ of CY dimension $3$, $\CC^\bullet(A,A)$ is an
$E_3$-algebra; the $E_3$-action is canonical up to a contractible
choice (associator torsor $\times$ $SO(3)/SO(2) = S^2$ of forgetful
maps $fE_4 \to E_3$).

\item[SC6 (CFG compatibility).] The three $E_3$-structures on
Hochschild-type complexes over a CY$_3$ (topological from real
$\R^3$-factorisation; holomorphic from complex $\C^3$-factorisation;
algebraic from CY$_3$-cyclic pairing) are compatible on each
$\C^3$-chart via Kontsevich--Tamarkin formality. Globally on a
compact CY$_3$, the compatibility requires the
$\Theta_{\mathrm{hCS} \to \mathrm{Hall}}$ comparison map (open).
\end{description}

\subsection*{Ghost theorems identified}

\begin{description}
\item[GT1.] ``$\CC^\bullet$ is an $E_2$-algebra'' without exhibiting
the brace-operad map. The ghost: one can \emph{always} name operadic
structure on Hochschild cochains, but only the brace model supplies
checked signs. Without chain-level signs, everything downstream
(cyclic BV, MC equation, shadow tower) is formal.

\item[GT2.] ``BV implies $E_3$'' without specifying which $E_3$.
The ghost: BV $=$ $fE_2$ on cohomology, not $fE_3$; the promotion
to $E_3$ needs the CY dimension as operadic shift data. Without this
the BV-to-$E_3$ step is a slogan.

\item[GT3.] ``Deligne at $n=3$'' conflating iterated Hochschild and
cyclic-CY$_3$-BV. The ghost: two different $E_3$-algebras on the
same complex, only identified after KT formality. A statement
``$E_3$ on $\CC^\bullet(\cC)$ for CY$_3$'' that does not name which
$E_3$ is a latent anti-pattern.

\item[GT4.] ``$E_3$ on Hochschild = $E_3$ from $\C^3$''. The ghost:
algebraic-$E_3$ from the cyclic pairing lives on $\CC^\bullet$; the
geometric $E_3$ from $\C^3$ lives on $\Obs_{\mathrm{hCS}}$. They
match by formality on each chart; globally they do not match
without a comparison map. Writing ``the $E_3$-structure'' as if
unique is type-confusion.

\item[GT5.] ``$R$-matrix is the $E_3$-bracket''. The ghost:
$R$-matrices at $d = 3$ come from the $E_2$-braiding on
$\Rep(D(Y^+))$, not from the $E_3$-bracket on $\CC^\bullet$. The
$E_3$-bracket governs deformation quantisation of $Y^+$; the
$R$-matrix emerges one Drinfeld-double step later. The full pipeline
is: $E_3$-MC on $\CC^\bullet \to E_1$-deformed CoHA $\to$ $E_2$-
braided double $\to$ $R$-matrix.
\end{description}

\subsection*{Paths for the chapter inscription}

\begin{enumerate}
\item Add a subsection ``Deligne at $n=2$ and at $n=3$ on cyclic
Hochschild cochains'' to the $E_3$-Hochschild section
(sec:e3-hochschild-deformation-theory) of
quantum\_chiral\_algebras.tex, inscribing the surviving core
SC1--SC5 as theorems with proofs, and SC6 as a remark.

\item Add the explicit chain-level $\Delta$ formula
\eqref{eq:delta-explicit} as a proposition with proof (Menichi
Prop.~2.1 transport).

\item Add a remark naming the three $E_3$-structures and the
compatibility chain, pointing to cy3\_chain\_level\_bridge.tex for
the global-on-compact-CY$_3$ open lemma.

\item Add a proposition naming the $\mathrm{GRT}_1$-torsor and the
$S^2$-choice of forgetful map $fE_4 \to E_3$, as the two sources of
non-uniqueness in the restricted algebraic $E_3$-action.

\item Remove any sentence that writes ``the Deligne conjecture at
$n=3$'' without naming which of A, B, C (iterated /
cyclic-$d=2$-BV / cyclic-$d=3$) is meant. Target chapter currently
has a few sentences that conflate; they are cleaned in the
inscription pass.
\end{enumerate}

\subsection*{Key chain-level identities (audit sheet)}

\textbf{Brace identity (Voronov arXiv:9908040 Thm.~3.3).}
For $x \in \CC^m, y_i \in \CC^{k_i}$ ($i = 1, \ldots, n$),
\[
 \{x; y_1, \ldots, y_n\}(a_1, \ldots, a_N)
 \;=\;
 \sum_{\mathrm{shuffles}} (-1)^\eta\,
 x(a_1, \ldots, y_1(\ldots), \ldots, y_n(\ldots), \ldots, a_N),
\]
$N = m + \sum k_i - n$, shuffled so each $y_j$ consumes $k_j$ consecutive arguments starting after $y_{j-1}$.

\textbf{Brace-associativity.}
\[
 \{\{x; y_1, \ldots, y_m\}; z_1, \ldots, z_n\}
 \;=\;
 \sum_{\mathrm{patterns}} (-1)^{\eta'}
 \{x; \ldots, \{y_j; z_?, \ldots\}, \ldots\}.
\]

\textbf{Cup product.} $(f \smile g)(a_1, \ldots, a_{m+n}) =
(-1)^{?} f(a_1, \ldots, a_m) \cdot g(a_{m+1}, \ldots, a_{m+n})$
(the arity-$2$ brace with no insertions); Koszul sign from moving
$g$ past $a_1, \ldots, a_m$.

\textbf{Gerstenhaber pre-Lie.}
$f \circ g = \{f; g\} = \sum_{i=1}^{m} (-1)^? f(a_1, \ldots, g(\ldots), \ldots, a_{m+n-1})$.

\textbf{Gerstenhaber bracket.}
$[f, g]_\Ger = f \circ g - (-1)^{(|f|-1)(|g|-1)} g \circ f$.

\textbf{BV identity.}
$[f, g]_\Ger = (-1)^{|f|}[\Delta(f \smile g) - \Delta(f) \smile g
 - (-1)^{|f|} f \smile \Delta(g)]$.

\textbf{Cyclic $\Delta$ (dimension $d$).}
\eqref{eq:delta-explicit}, with pairing-degree shift $-d$.

\textbf{$\Delta^2 = 0$ and $[b, \Delta] = 0$.} Explicit from
cyclic symmetry.

\textbf{Operadic zigzag.}
$\Br \otimes \Lambda[\Delta] \xrightarrow{\sim}
\mathrm{End}(\CC^\bullet(A,A))$ over
$C_\bullet(fE_{d+1};\Q)$.

\textbf{Higher Deligne iteration.}
$\CC^\bullet$ of $E_n$-algebra is $E_{n+1}$. Lurie HA Thm.~5.3.1.30:
for an $E_n$-algebra $A$ in spectra, the centre $Z(A)$ is $E_{n+1}$
and $\CC^\bullet(A,A) \simeq Z(A)$. At $n = 1$: $E_2$, the original
Deligne. At $n = 2$: $E_3$. At $n = 3$: $E_4$.

\subsection*{Side-by-side with CFG 2026}

CFG 2026 Appendix C constructs an $E_3$-action on observables of
topological Chern--Simons (with spectral parameter). Their derivation
uses:
\begin{itemize}
\item $\R^3$-factorisation and the little $3$-disks operad directly;
\item the deformation quantisation of the classical $E_3$ algebra
$C^\infty(\mathfrak{g}^* \otimes H^\bullet(\mathrm{pt}))$ of Wilson-line
functionals;
\item the ribbon-category output via factorisation-homology traces
over a solid torus.
\end{itemize}

The present agent's derivation uses:
\begin{itemize}
\item brace operad on Hochschild cochains, McClure--Smith;
\item cyclic $\Ainf$ structure on $A_X$ for a CY$_3$, Kontsevich--Soibelman;
\item BV descent Menichi--Tradler--Zeinalian;
\item formality bridge $\Br \otimes \Lambda[\Delta] \xrightarrow{\sim} fE_{d+1}$.
\end{itemize}

Common downstream: both produce an $E_3$-algebra that feeds into the
factorisation-homology machine. CFG 2026 is algebraically simpler
(no CY hypothesis) but topologically restricted ($\R^3$, no
holomorphic enhancement). The present construction is algebraically
richer (full CY$_3$-cyclic) and holomorphically compatible (matches
hCS on $\C^3$ charts via KT formality), but requires the cyclic
pairing input and the formality bridge.

Resolution of the ``CFG shortcut'' warning in
cy3\_chain\_level\_bridge.tex (warn:cy3-no-cfg-shortcut): CFG's
$E_3$ algebra at $d=3$ real is \emph{not} the same as the
CY$_3$-cyclic $E_3$; identifying them is the $\Theta_{\mathrm{hCS}
\to \mathrm{Hall}}$ open problem. The present agent's cycles
clarify \emph{which} $E_3$ on \emph{which} complex is produced by
\emph{which} mechanism, reducing the inter-operad confusion to a
single named morphism.

\subsection*{Status of the inscription}

The inscription target is
\texttt{chapters/theory/quantum\_chiral\_algebras.tex},
sec:e3-hochschild-deformation-theory (line 2346 onwards).
Material to add:
\begin{enumerate}
\item New subsection ``Deligne at $n=2$: brace operad on Hochschild
cochains''.
\item New subsection ``BV descent from the CY cyclic pairing''.
\item New subsection ``Deligne at $n=3$: three compatible
$E_3$-structures''.
\item Theorems: SC1, SC2, SC3, SC4, SC5 above.
\item Propositions: chain-level $\Delta$ formula; BV identity.
\item Remarks: GT1--GT5 cleaned of the ghost; CFG compatibility.
\end{enumerate}

The wave file concludes with the final audit sheet of chain-level
identities, ensuring that any future edit that reuses these formulas
does not reintroduce GT1--GT5.

\end{document}
